\chapter{Bayesian Inference, Graphical Models, and Learning under Uncertainty}

Deterministic predictions are often not enough. In many problems the main question is not only ``what is the most likely answer?'' but also ``how certain is the model, and how should that certainty change when new evidence arrives?'' Bayesian inference provides the formal language for that update.

\section{Bayes' rule}

For hypothesis $h$ and data $D$,
\[
  p(h \mid D)
  =
  \frac{p(D \mid h)\,p(h)}{p(D)}.
\]
The prior $p(h)$ encodes knowledge before the observation. The likelihood $p(D \mid h)$ scores how compatible the data are with that hypothesis. The posterior combines both.

\begin{keyequation}[Posterior update]
\[
p(h \mid D) \propto p(D \mid h)\, p(h)
\]
\tcblower
Bayesian inference is a competition between prior plausibility and data fit. The posterior reweights beliefs rather than replacing them from scratch.
\end{keyequation}

\section{Conditional independence}

Probabilistic models become useful only when they factorize. The simplest nontrivial notion is conditional independence:
\[
  X \perp Y \mid Z
  \quad \Longleftrightarrow \quad
  p(x,y \mid z) = p(x \mid z)\,p(y \mid z).
\]
This reduces both storage and inference cost. It is the reason graphical models scale at all.

\begin{aside}[Why independence matters]
Without independence assumptions, the joint distribution over many variables grows exponentially. Conditional independence is the mathematical device that turns an impossible joint table into a manageable model.
\end{aside}

\section{Bayesian networks}

A Bayesian network is a directed acyclic graph whose factorization is
\[
  p(x_1,\dots,x_n)
  =
  \prod_{i=1}^{n} p(x_i \mid \mathrm{pa}(x_i)),
\]
where $\mathrm{pa}(x_i)$ denotes the parents of node $x_i$.

\begin{figure}[t]
\centering
\begin{tikzpicture}[node distance=2.2cm, >=Latex]
  \node[draw, rounded corners] (b) {Burglary};
  \node[draw, rounded corners, right=of b] (e) {Earthquake};
  \node[draw, rounded corners, below right=1.5cm and 1.1cm of b] (a) {Alarm};
  \node[draw, rounded corners, below left=1.7cm and 0.8cm of a] (j) {John calls};
  \node[draw, rounded corners, below right=1.7cm and 0.8cm of a] (m) {Mary calls};
  \draw[->] (b) -- (a);
  \draw[->] (e) -- (a);
  \draw[->] (a) -- (j);
  \draw[->] (a) -- (m);
\end{tikzpicture}
\caption{A canonical Bayesian-network example. The graph represents factorization structure, not merely visual correlation.}
\end{figure}

The graph encodes more than a drawing. It says exactly how the joint distribution factorizes and which variables become irrelevant once certain others are known. The \emph{Markov blanket} of a node consists of its parents, children, and co-parents of its children; conditioned on that set, the node is independent of the rest of the graph.

\section{Bayesian prediction and model selection}

In Bayesian learning, model parameters are themselves uncertain. If $D$ is the dataset and $\wvec$ are parameters, then
\[
  p(\wvec \mid D)
  \propto
  p(D \mid \wvec)\, p(\wvec).
\]
Predictions average over that posterior:
\[
  p(y^\ast \mid \vec{x}^\ast, D)
  =
  \int p(y^\ast \mid \vec{x}^\ast, \wvec)\, p(\wvec \mid D)\,\dd \wvec.
\]

This formula is one of the cleanest statements in the course. It says prediction should average over plausible models, not commit too early to a single one.

\section{MAP estimation and regularization}

If the full posterior integral is too costly, one may approximate it by a single parameter choice,
\[
  \wvec_{\mathrm{MAP}}
  =
  \argmin_{\wvec}
  \left[
    -\log p(D \mid \wvec) - \log p(\wvec)
  \right].
\]
This reveals the familiar bridge between Bayesian and frequentist viewpoints: a log-prior acts like a regularizer. For example, a Gaussian prior on $\wvec$ yields an $\ell_2$ penalty.

\section{What Bayesian methods add}

The probabilistic block of the course contributes three durable ideas:
\begin{itemize}
  \item uncertainty should be represented explicitly rather than hidden,
  \item graphical structure is a computational tool,
  \item regularization can be interpreted as prior information.
\end{itemize}
